\section{Introduction}

In medical imaging, longitudinal data plays a central role in evaluating how anatomy changes over time. Clinicians rarely interpret a scan in isolation; instead, they often assess follow-up images with respect to a baseline, using the subject’s own imaging history as the reference point. This approach accounts for natural anatomical variability and enables a more personalized judgment of what constitutes as progression.

These trajectories are shaped by a range of subject-specific factors, including chronological age, pre-existing conditions, or genetic risk, that influence both the pattern and rate of change. An older subject, for example, may show substantial atrophy even in the absence of disease, while another may exhibit more abrupt degeneration due to pathological processes. In this context, the interpretation of a scan depends as much on where the subject started as on where they are now.

Yet despite its clinical value, longitudinal imaging data is often incomplete. Missed visits, dropouts, and inconsistent imaging protocols introduce sparsity across timepoints, making it difficult to reconstruct a coherent picture of disease progression. When available, complete trajectories offer the possibility of capturing gradual trends as well as detecting deviations from them—events that might otherwise be missed but could carry clinical significance. For data-driven models, such irregularity complicates the learning of coherent temporal patterns and introduces risk of bias in downstream tasks~\cite{ehrmissing2023, bell2023, little2009}.

Generative models offer a natural framework for scan imputation, and recent work in GANs and diffusion models has demonstrated strong results in image synthesis. Still, many of these methods operate directly in intensity space, generating voxel values anew. In high-dimensional medical images, this can lead to anatomically implausible results, especially when training data is limited. Some approaches synthesize 2D slices and reassemble them into volumes~\cite{yuan2024remind}, at the risk of discontinuities across planes. Others model deformation trajectories explicitly~\cite{bai2024noder}, but require separate models per subject and per timepoint, limiting generalization and scalability.

Additionally, in most longitudinal scans, most anatomical regions remain unchanged between timepoints. These areas offer limited insight into disease progression. It is the evolving regions—subtle, spatially localized, and shaped by each subject’s trajectory, that carry the signal of interest. From a clinical perspective, what matters is not the absolute appearance of a scan, but how it differs from where the subject began.

This suggests that modeling transformation, rather than appearance, might offer a more effective path forward. Anchoring all timepoints to a common baseline reflects how clinicians reason about progression, and provides a consistent frame of reference across individuals. In doing so, it also reduces inter-subject variability, which facilitates comparison.

This perspective is particularly relevant in conditions like Alzheimer’s disease, where two subjects may show comparable reductions in tissue volume at a given follow-up, though only one may be undergoing pathological decline. This distinction has significant clinical implications, especially in early-stage intervention and patient-specific monitoring. Modeling change relative to baseline helps distinguish such cases by restoring the temporal context. 

A well-established representation of anatomical change in neuroimaging is the stationary velocity field (SVF), widely used in deformable registration~\cite{ashburner2007fast, modat2012, fishbaugh2020}. SVFs define smooth, invertible transformations between images, offering a compact, geometrically grounded way to encode how anatomy deforms over time. Rather than synthesizing entire images, one estimates the transformation required to reach the desired anatomical state. This shifts the modeling focus from raw appearance to subject-specific progression, which is arguably the central object of interest in longitudinal analysis.

\vspace{0.5em}

In this paper, we introduce SynBADD (Longitudinal Medical Image Imputation via \textbf{Syn}thetic \textbf{B}idirectional \textbf{A}natomy-Guided \textbf{D}iffeomorphic \textbf{D}eformations), a generative framework for imputing missing 3D scans by predicting anatomically plausible deformation fields from the subject’s baseline. Rather than synthesizing voxel intensities, SynBADD estimates the diffeomorphic transformation that would align the baseline to the missing timepoint. This encourages consistency with anatomical structure, and constrains generation to plausible deformation trajectories. A transformer-based model learns to reason jointly over past and future changes, capturing nonlinear and individualized progression patterns. To address the scarcity of real longitudinal data, we pre-train a variational autoencoder on a large corpus of synthetically generated SVFs, learning a compact latent space of deformation dynamics. A refinement module then enhances the predicted field to preserve fine-grained anatomical detail.

\vspace{0.5em}

Our contributions are as follows:
\begin{itemize}
    \item We introduce a novel trajectory-conditioned deformation-based generative model that imputes missing intermediate scans by leveraging both past and future anatomical information across subjects.
    \item We propose a compact latent representation of SVFs learned from synthetic data, enabling effective learning under limited data by decoupling deformation modeling from raw intensity generation.
    \item We design a spatial refinement module that improves the fidelity of the predicted deformation, preserving subject-specific anatomical detail.
    \item We achieve state-of-the-art results on the ADNI dataset, improving SSIM by $2.39\%$ and PSNR by $10.64\%$, while maintaining anatomically coherent progression in synthesized scans.
\end{itemize}



% \section{Introduction}

% In medical imaging, longitudinal data plays a central role in evaluating how a subject changes over time. Clinicians rarely interpret a scan in isolation; instead, they often assess follow-up images with respect to a baseline, using the subject’s own history as the reference point. This approach accounts for natural anatomical variability and enables more personalized judgment of what constitutes progression.

% These trajectories are shaped by a range of subject-specific factors—chronological age, pre-existing conditions, genetic risk—that influence both the pattern and rate of change. An older subject, for example, may show substantial atrophy even in the absence of disease, while another may exhibit more abrupt degeneration due to pathological processes. In this context, the interpretation of a scan depends as much on where the subject started as on where they are now.

% Yet despite its clinical value, longitudinal imaging data is often incomplete. Missed visits, dropouts, and inconsistent imaging protocols introduce sparsity across timepoints, making it harder to reconstruct a coherent picture of disease progression. When available, complete trajectories offer the possibility of capturing gradual trends as well as detecting deviations from them—events that might otherwise be missed but could carry clinical significance.

% This setting presents modeling challenges distinct from those found in other temporal domains. Longitudinal imaging data evolves slowly, often over months or years, with much of the anatomical structure remaining stable between timepoints. Change, when it occurs, tends to be subtle, spatially localized, and subject-specific. This contrasts with modalities like fMRI or EEG, where high-frequency fluctuations dominate and population-level trends often guide analysis. As a result, techniques that assume homogeneous temporal dynamics—interpolation, regression, even autoregressive extrapolation—may fail to capture the nuanced, individualized nature of structural change.

% In most longitudinal scans, the majority of anatomical structure remains unchanged between timepoints. These stable regions, while expected, offer little insight into disease progression. What matters are the changes—often subtle and spatially localized—that signal underlying biological processes.

% This naturally raises the question of whether modeling full images is the most effective strategy. Since clinical assessments focus on change relative to baseline, it may be more appropriate to operate in a space where transformation—not appearance—is the primary signal. Anchoring all timepoints to the baseline provides a consistent frame of reference, reducing inter-subject variability and aligning with clinical reasoning.

% This perspective is particularly relevant in conditions like Alzheimer’s disease, where two subjects may show comparable reductions in tissue volume at a given follow-up, though only one may be undergoing pathological decline. Modeling change relative to baseline helps distinguish such cases by restoring the temporal context.

% In neuroimaging, a natural representation of such change is the stationary velocity field (SVF), widely used in registration tasks. SVFs describe smooth, invertible transformations between images and offer a compact way to encode anatomical change. Modeling in SVF space shifts attention from image synthesis to deformation estimation, isolating the meaningful signal and reducing the modeling burden.


% % In medical imaging, clinicians use longitudinal data to measure change in a subject with respect to a baseline. Each subject has a unique trajectory which is influenced by factors like their age which may show an already advanced degree of atrophy or existing disease conditions which can accelerate tissue degeneration. Due to the various factors which shape a scan, evaluating it at a single timepoint is not enough and we need the surrounding context.

% % These longitudinal datasets are often incomplete or sparse due to various reason including but not limited to patient drop out, missed visit and protocol variability. Completeness of these datasets help in providing a fuller picture of disease evolution and can help in noticing when there is a sudden fluctuation as it may indicate pathology.

% % Longitudinal medical data is unique from other temporal data as it pertains to a single subject over a long period of time which may range from a few months to several years. Most of the subject remains the same but there are subtle changes which are influenced by the above mentioned conditions. The nature of processes which affect longitudinal images are thus different from other types of temporal data, like fMRI or EEG etc.

% % Subjects have a unique trajectory which is influenced by a velocity term which depends upon their age and an acceleration term which may be affected by medical conditions~\cite{}. To effectively model such trajectories, we need to move beyond frameworks like interpolation, regression or extrapolation which are more suited to other temporal modalities operating over shorter windows where the data evolution is governed by short term forces.

% % Due to each subject having such a unique trajectory, the structural data becomes difficult to interpret in a stand alone manner. what clinicians do to handle this inter subject variability and ellicit common trends for diagnosis, is to look at the changes in the context of the subject, as a difference to the baseline.

% % Operating in change space provides contextualizes the data per subject, which makes it free of inter-subject variations like starting age, disease label etc.,  allowing us to compare trajectories across subjects. In other words, when we look at raw images, an older subject might present greater atrophy despite being healthy due to normal aging. However, without the context of the baseline scan present, comparing the month 18 visit for an older subject with a younger subject, might cause the older subject's case to appear as a pathology. Thus, operating in change space helps to contextualize each timepoint/scan or data point. 
% % ---

% % We attempt to take a look at longitudinal image imputation in the context of AD.


% % --- Importance of clinical trajectory and it's relevance. Motivate subtely to AD.
% % This is achieved through longitudinal imaging. but it is not very feasible. And hence the need for it.
% % --- Change is clinically relevant.  Clincians use longitudinal data to see change over time. Particularly in the region of interest. In AD, give examples. (structural imaging in not the gold standard in AD; only use AD as an example; AD provides the tech challenge but it's not the perfect use case).
% % --- Why imputation in longitudinal imaging is interesting -- basically, how different is it from other types of imputation. (why is it worth working on). -- subtle change while preserving subject identity. (what should change and what shouldn't change)
% % --- Challenges: why is temporal trajectory modelling important for this task and challenging (and how is it different from just interpolation and regression). 
% % --- Gently put the idea into their heads (why we operate in change space; make them agree that change space is better for longitudinal imaging) - Why change? Clincians use longitudinal data not for the trend but to see change relative to that specific subject. How do measure change -- SVF (known and well used in medical image analysis).

% % ---

% % Understanding and tracking how an individual’s anatomy changes over time is essential for both clinical research and precision medicine. Unlike cross-sectional imaging, which captures a static snapshot at a single timepoint, longitudinal imaging allows direct observation of anatomical evolution within a subject. This enables the study of normal aging as well as disease progression, offering valuable insight into how conditions unfold. A major strength of longitudinal imaging is its ability to disentangle within-subject temporal changes from between-subject variability, enabling more accurate characterization of disease trajectories, personalized prognosis, and treatment planning. As a result, longitudinal study designs have become central to population-level studies and progressive disease monitoring.\\

% % In practice, however, longitudinal imaging datasets are often small and suffer from missing or sparsely sampled timepoints. Acquiring high-resolution 3D scans over long periods is expensive, logistically difficult, and often infeasible in routine clinical settings. Even in large-scale cohorts like the Alzheimer’s Disease Neuroimaging Initiative (ADNI), intermediate timepoints are frequently missing due to subject dropout, missed visits, protocol variability, cost constraints, or loss to follow-up. These missing scans may coincide with clinically significant transitions, such as the onset of rapid decline or transient treatment effects, that are not evident from baseline or follow-up scans alone. Their absence makes it difficult to infer what may have occurred between two observations, especially in retrospective studies where reacquisition is impossible. This limits the ability to trace disease trajectories over time, which is often necessary for clinical audits, interpretation of findings, and treatment planning. From a machine learning standpoint, missing timepoints disrupt temporal continuity, hinder effective modeling of disease trajectories, and introduce bias in outcome prediction.\\

% % Much of prior work in longitudinal image synthesis focuses on predicting future scans from earlier timepoints. In contrast, imputing missing intermediate scans presents a distinct set of challenges. Unlike extrapolation, which projects anatomy forward in time, intermediate imputation requires reasoning over the entire temporal window, integrating both past and future scans to generate anatomically and temporally coherent estimates. This is particularly important in disease populations, where anatomical changes are often nonlinear, subtle, and highly variable across individuals. In such cases, future scans may provide crucial evidence on what must have occurred earlier—for instance, the presence of rapid atrophy in the future may imply earlier degeneration that is not apparent from the past alone. Models that ignore this bidirectional temporal context or rely solely on forward prediction risk generating scans that are anatomically plausible in isolation, yet inconsistent with the subject's true trajectory.\\

% % Recent advances in generative models, particularly GANs and diffusion models, have achieved remarkable performance in photorealistic synthesis and cross-domain image translation. However, their application to longitudinal medical imaging remains limited. Most existing approaches generate directly in the intensity space, synthesizing voxel-wise values from scratch at each timepoint. In high resolution 3D imaging settings with limited training data, such direct synthesis is prone to producing unrealistic or anatomically inconsistent outputs. The challenge is even more pronounced when only a few scans per subject are available. Moreover, intensity-based synthesis lacks explicit structural priors, making it difficult to preserve subject identity or ensure consistency across time. For example, ReMiND generates 2D slices independently using a diffusion model and then stacks them into volumes, but lacks volumetric reasoning and temporal awareness, often resulting in broken anatomical continuity. Separately, NODER models deformation trajectories using neural ODEs but trains a separate model per subject, limiting population-level generalization and robustness under sparse sampling.\\

% % In this work, we introduce SynBADD, a structured generative framework for imputing missing intermediate scans using bidirectional anatomy-guided diffeomorphic deformations. Instead of synthesizing full intensity volumes, SynBADD models anatomical evolution via deformation fields that warp a reference scan to the missing timepoint. This formulation constrains synthesis to a plausible manifold of anatomical deformations, preserving spatial structure and subject identity. Our method leverages a transformer-based architecture to jointly reason over deformation trajectories from both past and future scans, capturing subject-specific aging and disease progression patterns for temporally aware synthesis. To handle limited real longitudinal data, we pretrain a variational autoencoder on a large corpus of synthetic stationary velocity fields (SVFs), learning a compact latent representation of anatomical deformations. The generated coarse deformation estimate is further refined using a spatial decoder to recover fine-grained anatomical detail. This analysis-by-synthesis strategy supports anatomically coherent, temporally constrained, and data-efficient scan imputation.\\

% % Our key contributions are as follows:
% % \begin{itemize}
% %     \item We introduce SynBADD, a trajectory-conditioned, deformation-based generative framework for imputing missing intermediate scans in longitudinal medical imaging via bidirectional, anatomy-guided diffeomorphic deformations. To the best of our knowledge, it is the first method to use a transformer-based trajectory model over latent deformation fields for intermediate scan imputation.
% %     \item We propose a compact latent representation of stationary velocity fields (SVFs), learned via a variational autoencoder pretrained on synthetically augmented deformations. This enables deformation modeling decoupled from direct intensity synthesis, supporting data-efficient learning under limited real-world longitudinal data.
% %     \item We develop a bidirectional transformer with learnable temporal and time directionality embeddings to reason jointly over past and future anatomical changes, capturing subject-specific, temporally consistent deformation patterns.
% %     \item We design a spatial refinement module that enhances the fidelity of the predicted deformation fields and preserves detailed anatomical structure in the imputed scans.
% %     \item We demonstrate that SynBADD outperforms state-of-the-art baselines on the ADNI dataset, achieving a $XX\%$ improvement in SSIM and a $XX\%$ gain in PSNR and $XX\%$ higher Dice scores in regions affected by neurodegeneration, while more faithfully capturing anatomically consistent changes at intermediate timepoints.
% % \end{itemize}

% % The remainder of this paper is organized as follows. Section 2 reviews relevant literature, beginning with challenges of missing data in longitudinal medical imaging and followed by prior approaches for future scan prediction, generative modeling, and deformation-based synthesis. We conclude the section with a discussion of Alzheimer’s disease as a motivating application and the role of longitudinal imaging in neurodegenerative disorders.  Section 3 introduces the proposed SynBADD framework, detailing its latent deformation modeling, bidirectional temporal inference, and anatomically grounded synthesis pipeline. Section 4 describes the experimental setup, including dataset preprocessing, training protocols, and evaluation metrics. Section 5 presents both quantitative and qualitative results, compares SynBADD against state-of-the-art baselines, and discusses key findings. Finally, Section 6 concludes the paper and outlines potential directions for future work.